% chapters/appendix-ch15-sql.tex
\chapter{SQL Appendix: Chapter 15 (Design Trade-offs)}

\section{DDL (denormalized read model)}
\begin{verbatim}
CREATE TABLE loan (
  loan_id INT PRIMARY KEY AUTO_INCREMENT,
  member_id INT NOT NULL,
  checkout_date DATE NOT NULL,
  due_date DATE NOT NULL,
  return_date DATE NULL
) ENGINE=InnoDB;

-- Denormalized read model (trade-off example)
CREATE TABLE loan_read_model (
  loan_id INT PRIMARY KEY,
  member_id INT NOT NULL,
  due_date DATE NOT NULL,
  overdue_flag TINYINT(1) NOT NULL DEFAULT 0
) ENGINE=InnoDB;
\end{verbatim}

\section{Sample data}
\begin{verbatim}
INSERT INTO loan (member_id, checkout_date, due_date, return_date)
VALUES (1, '2026-01-10', '2026-01-24', NULL);

INSERT INTO loan_read_model (loan_id, member_id, due_date, overdue_flag)
VALUES (1, 1, '2026-01-24', 0);
\end{verbatim}

\section{Example queries}
\begin{verbatim}
-- Q1: Read model usage
SELECT loan_id, member_id, overdue_flag
FROM loan_read_model
WHERE overdue_flag = 1;

-- Q2: Recompute overdue flag (derived truth)
SELECT loan_id,
       CASE WHEN return_date IS NULL AND due_date < CURRENT_DATE()
            THEN 1 ELSE 0 END AS derived_overdue
FROM loan;
\end{verbatim}
